\section{From research gaps to the problem statement}

\subsection{Where the project started}

At mid-semester the project was a small prototype: 5 sensors, randomly generated data, a max-age baseline and a
heuristic that discounted correlated sensors (we called it ``v4''). Our faculty asked for two things: \textbf{use many
more sensors} and \textbf{use real data}, and show a live demonstration of how the scheduling happens. While
rebuilding, we also found a bug in the prototype: the correlation graph was estimated on one random sensor layout
(seed 42) but evaluated on a different one (seed 123), so the ``correlation-aware'' numbers were not meaningful. That
pushed us to rebuild the project properly and to read the literature carefully to position it.

\subsection{What the literature does (2021--2026)}

We grouped recent work into three streams. Full references are in Section~\ref{sec:refs}.

\textbf{Stream A: AoI with correlated sources (analytical).}
\begin{itemize}
  \item Tripathi \& Modiano \cite{1} (ACM MobiHoc 2022; IEEE/ACM ToN 2024): randomised and max-weight policies for
        weighted AoI when an update from one source can refresh others. Within a factor 2 of optimal. Correlation is an
        abstract probability, the objective is AoI, not estimation error.
  \item Ramakanth, Tripathi \& Modiano \cite{2} (IEEE INFOCOM 2024; IEEE TMC 2025), \textbf{our base paper}: correlated
        sources as a Wiener process, Kalman-filter estimation, one source observed per slot. Shows the error of any
        policy is bounded by a weighted sum of AoI, so AoI-based scheduling that ignores correlation is within a
        \emph{constant factor} of optimal.
  \item Erbayat \emph{et al.}\ \cite{3,4} (ACM MobiHoc 2024, IEEE ISIT 2025): queueing analysis of AoI and state error
        when one sensor's update covers several processes; preemption strategies.
  \item Fidler \emph{et al.}\ \cite{5} (IEEE Trans.\ Commun.\ 2026): 2D-AoI, which converts spatial distance into an
        equivalent age. A metric, not a scheduler.
\end{itemize}

\textbf{Stream B: learning-based AoI and estimation scheduling.}
\begin{itemize}
  \item Ceran, G\"und\"uz \& Gy\"orgy \cite{6} (IEEE JSAC 2021): RL for AoI with retransmissions (HARQ); sources
        independent.
  \item Farag \emph{et al.}\ \cite{7} (IEEE GCAIoT 2021): actor-critic DRL for average AoI plus the probability of
        exceeding an AoI threshold (threshold met only in probability).
  \item Chen \emph{et al.}\ \cite{8} (IEEE ICC 2023): structure-enhanced DRL for remote-estimation MSE; 10--25\% lower
        MSE than standard DRL. Needs the optimal policy's structure derived by hand first.
  \item Reyhan \& Coleri \cite{9} (IEEE Trans.\ Commun.\ 2025): safe DRL with peak-AoI \emph{violation probability}
        constraints; optimises power.
\end{itemize}

\textbf{Stream C: graph neural networks for wireless scheduling.}
\begin{itemize}
  \item Shen \emph{et al.}\ \cite{10} (IEEE JSAC 2021): message-passing GNNs for radio resource management; generalise
        to larger networks.
  \item Zhao \emph{et al.}\ \cite{11} (IEEE TWC 2023): GCN-guided link scheduling (max-weight independent set).
  \item Wang \emph{et al.}\ \cite{12} (IEEE/IFIP WiOpt 2024): MPNN that optimises an age-aware randomised policy.
  \item Wang \emph{et al.}\ \cite{13} (IEEE TWC 2026): federated GNN + multi-agent RL for AoI-aware task offloading.
\end{itemize}

\subsection{The research gaps}

Reading these side by side, five gaps stood out:

\begin{description}[style=nextline,leftmargin=1.2cm]
  \item[G1 -- Correlation is hand-modelled, not learned.] Stream A assumes a known correlation model (probabilities,
        Wiener covariance, Gaussian kernels). Real deployments do not come with one.
  \item[G2 -- The objective is freshness, not knowledge.] Most work minimises AoI. AoI is a proxy; what the sink
        actually needs is a low estimation error for \emph{all} sensors.
  \item[G3 -- Learned schedulers use the wrong graph.] GNN schedulers (Stream C) build graphs from interference or
        road topology, i.e.\ who \emph{conflicts} with whom, not whose \emph{data is redundant}.
  \item[G4 -- Freshness is soft.] Learned methods treat AoI limits as penalties or probabilistic constraints
        \cite{7,9}; a sensor can still be neglected.
  \item[G5 -- Little evidence on real multi-sensor data.] Almost all results are on simulations; few use real
        recordings with a strict separation between training and test periods. The base paper also polls only
        one source per slot and leaves a constant-factor gap open.
\end{description}

\subsection{How the gaps became our problem statement}

Each gap fixed one part of the problem statement:
\begin{center}
\begin{tabularx}{\textwidth}{@{}lX@{}}
\toprule
\textbf{Gap} & \textbf{What it added to our problem statement} \\
\midrule
G1 & The correlation structure must be \emph{learned from historical data}. \\
G2 & The objective is the sink's \emph{estimation error} over all sensors, not AoI. \\
G3 & The scheduler should be a GNN over a \emph{data-correlation} graph (and the sink's live uncertainty). \\
G4 & Every sensor's AoI must stay below $F$ \emph{by construction}, for every policy. \\
G5 & Evaluate on \emph{real} datasets with a chronological train/test split, with $C>1$ polls per slot. \\
\bottomrule
\end{tabularx}
\end{center}

\begin{keybox}[Final problem statement]
Given $M$ sensors whose readings are correlated in space and time, and a channel that carries at most $C$
transmissions per slot, design a scheduler that chooses which sensors transmit in each slot so as to
\textbf{minimise the sink's estimation error over all sensors}, using a correlation structure \textbf{learned from
data}, while \textbf{guaranteeing that no sensor goes more than $F$ slots without reporting}, and validate it on
\textbf{real sensor recordings} that were not used for training.
\end{keybox}

\subsection{The base paper and how we extend it}

\textbf{Base paper:} R.~V.~Ramakanth, V.~Tripathi, E.~Modiano, ``Monitoring Correlated Sources: AoI-based Scheduling
is Nearly Optimal,'' IEEE INFOCOM 2024 \cite{2}.

\begin{center}
\begin{tabularx}{\textwidth}{@{}lXX@{}}
\toprule
& \textbf{Base paper} & \textbf{Our project} \\
\midrule
Source model & Discrete-time Wiener process, known covariance & Stationary AR(1), covariance \emph{fitted from data} \\
Estimator & Kalman filter & Kalman filter (same idea) \\
Polls per slot & 1 & $C\ge1$ (6 and 40 on real data) \\
Policy & AoI-based, correlation ignored & Learned GNN, correlation-aware \\
Freshness & Not a constraint & Hard bound $F$ \\
Evaluation & Theory + simulation & Simulation + two real datasets, statistical tests \\
\bottomrule
\end{tabularx}
\end{center}

The base paper's AoI-based policy is exactly our \textbf{Max-AoI (B1)} baseline (in its unweighted form). The paper
proves it is within a constant factor of optimal; our results show that in practice the gap is real: our learned
scheduler lowers error by 20\% on Intel Lab against it.

\begin{saybox}
``Our base paper proves that simple AoI scheduling is within a constant factor of optimal when you monitor
correlated sources with a Kalman filter. That constant factor can still be large. We keep the same Kalman framework,
but learn the correlations from real data, poll several sensors per slot, guarantee freshness, and train a GNN to
close that gap. On the Intel lab data we close it by 20 percent.''
\end{saybox}

\subsection{Objectives and how each was achieved}

\begin{center}
\begin{tabularx}{\textwidth}{@{}p{0.6cm}p{5.2cm}X@{}}
\toprule
\textbf{\#} & \textbf{Objective} & \textbf{How we achieved it} \\
\midrule
O1 & Model the sink's knowledge with a Kalman filter on a correlation graph learned from data
   & \texttt{fit\_model} fits each sensor's dynamics and the innovation covariance on the training period;
     the graph keeps innovation correlations $\ge 0.3$; \texttt{KalmanSink} tracks estimate and covariance. (5.2--5.5) \\
O2 & Design a GNN scheduler trained end-to-end with PPO on estimation error
   & Three-layer GNN over two graphs + residual pointer decoder, initialised as Kalman-greedy, trained with PPO
     on $-\log$(error). (5.8--5.9) \\
O3 & Guarantee freshness by construction
   & Freshness guard forces overdue sensors before any policy acts; the environment asserts it every slot;
     unit tests check every baseline. (5.6) \\
O4 & Validate on real data with no leakage
   & Intel Lab (45 sensors) and PEMS-BAY (325); chronological split; daily profile and model fitted on training
     data only; a test checks no test data leaks. (6.5) \\
O5 & Compare fairly with significance tests and an ablation
   & Five baselines + oracle + ablation, identical readings for all, 95\% CIs, win rates, Wilcoxon tests. (5.10, 8) \\
O6 & Demonstrate live
   & Web demo replaying a real test day slot by slot with three schedulers side by side. (6.8) \\
\bottomrule
\end{tabularx}
\end{center}
